\documentclass[12pt,a4paper]{article}
\usepackage{expediente}
\begin{document}
% =============================================================================
% FICHA DE DATOS DEL EXPEDIENTE --- actualizado 20 ago 2026
% =============================================================================

% --- VENDEDOR / TITULAR REGISTRAL -------------------------------------------
\newcommand{\VendedorNombres}{RICHAR DARWIN}
\newcommand{\VendedorApellidos}{CUTIRE CONDORI}
\newcommand{\VendedorNombreCompleto}{RICHAR DARWIN CUTIRE CONDORI}
\newcommand{\VendedorDNI}{42610690}
\newcommand{\VendedorEstadoCivil}{soltero, sin unión de hecho vigente}
\newcommand{\VendedorOcupacion}{\faltante{[ocupación de Richar]}}
\newcommand{\VendedorDomicilio}{\faltante{[domicilio actual de Richar]}}
\newcommand{\VendedorTelefono}{\faltante{[teléfono de Richar]}}
\newcommand{\VendedorCorreo}{\faltante{[correo]}}
\newcommand{\VendedorNacionalidad}{peruano}
\newcommand{\VendedorFechaNac}{\faltante{[fecha de nacimiento]}}
\newcommand{\VendedorDNIEmision}{\faltante{[fecha de emisión del DNI de Richar --- la pide SUNARP]}}

% --- CONVIVIENTE: NO CORRESPONDE -------------------------------------------
\newcommand{\ConvivienteNombre}{NO CORRESPONDE}
\newcommand{\ConvivienteDNI}{NO CORRESPONDE}
\newcommand{\ConvivienteDesde}{NO CORRESPONDE}
\newcommand{\ConvivienteInscrita}{no existe unión de hecho vigente}
\newcommand{\ConvivienteDomicilio}{NO CORRESPONDE}

% --- COMPRADORES (copropiedad en partes iguales) ---------------------------
\newcommand{\CompradorUnoNombre}{NICO ALVARO CUTIRE ARCE}
\newcommand{\CompradorUnoDNI}{48325017}
\newcommand{\CompradorUnoEmision}{01 de enero de 2024}
\newcommand{\CompradorUnoEstadoCivil}{soltero}
\newcommand{\CompradorUnoOcupacion}{ingeniero informático}
\newcommand{\CompradorUnoDomicilio}{Jr.\ Wiracocha, Sicuani, provincia de Canchis, Cusco, Perú}
\newcommand{\CompradorUnoTelefono}{984715108}

\newcommand{\CompradorDosNombre}{LUCY CUTIRE ARCE}
\newcommand{\CompradorDosDNI}{47478852}
\newcommand{\CompradorDosEmision}{09 de enero de 2026}
\newcommand{\CompradorDosEstadoCivil}{soltera}
\newcommand{\CompradorDosOcupacion}{contador}
\newcommand{\CompradorDosDomicilio}{Comunidad Hanocca, distrito de Layo, provincia de Canas, Cusco, Perú}
\newcommand{\CompradorDosTelefono}{953440875}

\newcommand{\CompradorNombres}{NICO ALVARO y LUCY}
\newcommand{\CompradorApellidos}{CUTIRE ARCE}
\newcommand{\CompradorNombreCompleto}{NICO ALVARO CUTIRE ARCE y LUCY CUTIRE ARCE}
\newcommand{\CompradorDNI}{48325017 y 47478852}
\newcommand{\CompradorEstadoCivil}{solteros}
\newcommand{\CompradorConyuge}{NO CORRESPONDE}
\newcommand{\CompradorConyugeDNI}{NO CORRESPONDE}
\newcommand{\CompradorOcupacion}{ingeniero informático y contador, respectivamente}
\newcommand{\CompradorDomicilio}{Sicuani (Canchis) y Layo (Canas), Cusco}
\newcommand{\CompradorTelefono}{984715108 / 953440875}
\newcommand{\CompradorCorreo}{\faltante{[correo de contacto]}}
\newcommand{\CompradorNacionalidad}{peruanos}
\newcommand{\PorcentajeCopropiedad}{50\,\% cada uno}

% --- INTERMEDIARIO / OCUPANTE DEL REMANENTE --------------------------------
\newcommand{\IntermediarioNombre}{LUCIO GIL CUTIRE PARI}
\newcommand{\IntermediarioDNI}{42467850}
\newcommand{\IntermediarioEstadoCivil}{\faltante{[estado civil de Lucio]}}
\newcommand{\IntermediarioOcupacion}{\faltante{[ocupación]}}
\newcommand{\IntermediarioDomicilio}{\faltante{[domicilio de Lucio]}}
\newcommand{\IntermediarioTelefono}{\faltante{[teléfono de Lucio]}}
\newcommand{\IntermediarioParentesco}{familiar (apellido Cutire)}

% --- PREDIO MATRIZ ----------------------------------------------------------
\newcommand{\PartidaMatriz}{\faltante{[N.º de partida electrónica SUNARP --- pendiente]}}
\newcommand{\OficinaRegistral}{Oficina Registral de Puerto Maldonado}
\newcommand{\AreaMatriz}{\faltante{[área total inscrita; las partes refieren del orden de 100 ha]}}
\newcommand{\UnidadCatastral}{\faltante{[código catastral / UC, si existe]}}
\newcommand{\NombreFundo}{\faltante{[nombre del fundo, si tiene]}}
\newcommand{\KmCarretera}{Interoceánica Sur, al este del Puente Jayave (el puente queda al oeste del predio, sobre el río; no debajo del lote)}

% --- PREDIO A INDEPENDIZAR --------------------------------------------------
\newcommand{\AreaIndependizarHa}{10{,}00}
\newcommand{\AreaIndependizarMetros}{100 000{,}00}
\newcommand{\PerimetroM}{2 200{,}58}
\newcommand{\FrenteM}{100{,}29}
\newcommand{\FondoM}{1 000{,}00}

\newcommand{\PUnoLat}{$-12{,}912452$}
\newcommand{\PUnoLon}{$-70{,}170058$}
\newcommand{\PDosLat}{$-12{,}912497$}
\newcommand{\PDosLon}{$-70{,}170981$}
\newcommand{\PTresLat}{$-12{,}903469$}
\newcommand{\PTresLon}{$-70{,}171438$}
\newcommand{\PCuatroLat}{$-12{,}903424$}
\newcommand{\PCuatroLon}{$-70{,}170515$}

\newcommand{\PUnoE}{373 065{,}00}
\newcommand{\PUnoN}{8 572 256{,}00}
\newcommand{\PDosE}{372 965{,}00}
\newcommand{\PDosN}{8 572 251{,}00}
\newcommand{\PTresE}{372 911{,}00}
\newcommand{\PTresN}{8 573 249{,}00}
\newcommand{\PCuatroE}{373 011{,}00}
\newcommand{\PCuatroN}{8 573 255{,}00}

\newcommand{\FechaPago}{15 de agosto de 2018}
\newcommand{\MontoPago}{S/\ 20 000{,}00}
\newcommand{\MontoPagoLetras}{VEINTE MIL Y 00/100 SOLES}
\newcommand{\FormaPagoHistorica}{dinero en efectivo, sin medio de pago del sistema financiero}

% Nombres genéricos --- sustituir por los reales al ir a notaría
\newcommand{\NotarioNombre}{Dr.\ JUAN CARLOS RÍOS SALINAS (nombre genérico a sustituir)}
\newcommand{\NotariaCiudad}{Puerto Maldonado}
\newcommand{\ColegiaturaAbogado}{CAL Madre de Dios N.º 0000 (genérico)}
\newcommand{\AbogadoNombre}{Abog.\ ANA MARÍA QUISPE HUAMÁN (nombre genérico a sustituir)}
\newcommand{\IngenieroNombre}{Ing.\ PEDRO ANTONIO MAMANI CONDORI (nombre genérico a sustituir)}
\newcommand{\IngenieroCIP}{CIP 000000 (genérico)}
\newcommand{\Municipalidad}{Municipalidad Distrital de Inambari}
\newcommand{\Provincia}{Tambopata}
\newcommand{\Departamento}{Madre de Dios}
\newcommand{\Distrito}{Inambari}

\newcommand{\LugarOtorgamiento}{Puerto Maldonado}
\newcommand{\DiaOtorgamiento}{\faltante{[día]}}
\newcommand{\MesOtorgamiento}{\faltante{[mes]}}
\newcommand{\AnioOtorgamiento}{2026}

\newcommand{\TestigoUno}{\faltante{[testigo 1: nombres, DNI, domicilio]}}
\newcommand{\TestigoDos}{\faltante{[testigo 2: nombres, DNI, domicilio]}}

\InicioDocumento{Guía del proceso y orden del expediente}{DOC-02}

\noindent El expediente sigue el orden en que, en una casuística cooperante como la presente, suelen tramitarse los actos en Madre de Dios. Cada pieza tiene un PDF independiente en \texttt{Final Docs List}.

\Clausula{Hechos que este expediente toma como premisa}
El predio matriz está inscrito en SUNARP a nombre de \VendedorNombreCompleto, DNI \VendedorDNI. Las 10,00 ha son un recorte de un área mayor (aproximadamente la décima parte). En o alrededor del \FechaPago\ el comprador entregó \MontoPago\ en efectivo, documentado en hoja simple titulada compraventa, sin notario ni abogado. El dinero se canalizó a través de un intermediario que pretendía adquirir el área mayor y que hoy ocupa el predio. El comprador \textbf{no tomó posesión ni cultiva}. El titular vive, es soltero, tiene conviviente e hijos, puede acudir a notaría y está de acuerdo en firmar \textbf{directamente} con el comprador. Los tres (titular, intermediario-ocupante y comprador) están en condiciones de avanzar. El objetivo es inscribir las 10 ha a nombre del comprador.

\Clausula{Por qué no se ``legaliza el efectivo'' de 2018}
El pago ya ocurrió fuera del sistema financiero. En 2018 la transferencia de inmuebles ya exigía medio de pago desde 3 UIT; \MontoPago\ superaba ese umbral. En 2026 no se repara eso depositando la misma suma. La vía limpia es \textbf{reconocer el pago histórico} en la escritura y formalizar el dominio y el área. Tampoco se recomienda recaracterizar el acto como donación: el papel dice compraventa, hubo precio y el vendedor tiene herederos forzosos.

\Clausula{Secuencia operativa recomendada}
\begin{enumerate}[leftmargin=1.15cm, itemsep=0.28em]
\item \textbf{Due diligence registral.} Búsqueda de índices y obtención de copia literal y CRI (DOC-04 y DOC-05). Verificar que el polígono inscrito \textbf{contiene} los vértices P1 a P4.
\item \textbf{Cruce de riesgos territoriales.} Superposición con concesiones mineras (GEOCATMIN), patrimonio forestal, faja de derecho de vía de la Interoceánica Sur y catastro municipal.
\item \textbf{Levantamiento topográfico.} Profesional habilitado reemplaza el plano de trabajo (DOC-08 a DOC-11) por plano UTM, memoria descriptiva y, de ser exigido, visación / certificado catastral o negativo de zona catastrada.
\item \textbf{Saneamiento predial.} Constancia de no adeudo o convenio de pago en la \Municipalidad\ (DOC-06).
\item \textbf{Minutas.} Independización (DOC-13) y compraventa (DOC-14), con reconocimiento del pago 2018 (DOC-15), cesión del intermediario (DOC-16), consentimiento de conviviente (DOC-17) y acta de entrega de posesión (DOC-18).
\item \textbf{Alcabala.} Liquidación y pago en la municipalidad (DOC-07) antes o al formalizar, según práctica de la notaría.
\item \textbf{Escritura pública.} El notario eleva las minutas (DOC-19).
\item \textbf{Inscripción SUNARP.} Parte notarial vía SID-SUNARP (DOC-20). Apertura de partida independizada a nombre del comprador.
\item \textbf{Post-cierre.} Actualizar predial a nombre del comprador; archivar testimonio y copias (DOC-21).
\end{enumerate}

\Clausula{Mapa de documentos de este juego}
{\small\renewcommand{\arraystretch}{1.18}
\noindent\begin{tabularx}{\textwidth}{>{\bfseries}c >{\raggedright\arraybackslash}p{4.6cm} X}
\toprule
PDF & Pieza & Función \\
\midrule
01 & Carátula & Identifica el expediente. \\
02 & Esta guía & Orden del proceso. \\
03 & Cuestionario & Datos que faltan. \\
04--05 & Solicitudes SUNARP & Partida, literal y CRI. \\
06--07 & Solicitudes municipales & Predial y alcabala. \\
24 & Recibo trato directo 20k & Efectivo Richar--Nico/Lucy; Lucio acepta. No sustituye escritura. \\
08--11 & Técnico de trabajo & Vértices, memorias y plano. \\
12 & DJ estado civil & Vendedor y conviviente. \\
13 & Minuta de independización & Recorte de 10 ha y remanente. \\
14 & Minuta de compraventa & Título oneroso Cutire $\rightarrow$ comprador. \\
15 & Reconocimiento de pago 2018 & Prueba del precio histórico. \\
16 & Cesión y desistimiento & Cierra la cadena del intermediario. \\
17 & Consentimiento de conviviente & Reduce riesgo de impugnación. \\
18 & Acta de posesión y linderación & Entrega material del área. \\
19 & Proyecto de escritura & Cómo quedaría el instrumento. \\
20 & Solicitud de inscripción & Formulario tipo al Registro de Predios. \\
21 & Checklist de cierre & Control de firmas y requisitos. \\
25 & Casuística remanente Lucio & Compraventa vs donación de las $\sim$90 ha. \\
26 & Lucio a S/\ 45\,000 & Proceso, alcabala e IR de Richar. \\
\bottomrule
\end{tabularx}}

\Clausula{Qué no está en este juego (lo emite un tercero)}
Copia literal, CRI, certificado de búsqueda catastral, plano visado por verificador, constancia municipal de no adeudo, boleta de alcabala, DNI de las partes, parte notarial SID-SUNARP y asiento de inscripción. Este expediente deja \textbf{la solicitud} y el \textbf{espacio} para archivar cada original.

\Clausula{Advertencia sobre posesión}
Mientras el intermediario ocupe y el comprador no reciba la posesión, el dominio inscrito puede coexistir con un conflicto posesorio. Por eso el DOC-16 y el DOC-18 no son accesorios: son condición práctica del cierre.

\LugarFecha
\noindent En señal de haber leído esta guía de trabajo, las partes pueden rubricar:

\vspace{0.6em}
\TresFirmas{EL TITULAR / VENDEDOR}{\VendedorNombreCompleto\ --- DNI \VendedorDNI}{EL COMPRADOR}{\CompradorNombreCompleto\ --- DNI \CompradorDNI}{EL INTERMEDIARIO}{\IntermediarioNombre\ --- DNI \IntermediarioDNI}

\end{document}
